\section*{Chapter \ref{ch:b3:clt} --- Characteristic
Functions and the Central Limit Theorem}

\begin{solution}{exo:b3:clt:1}
Uniform on $\intcc{-1}1$: $\varphi(\xi) =
\frac12\int_{-1}^1\eu^{\iu\xi x}\dd x =
\frac{\sin\xi}{\xi}$ (equal to $1$ at $\xi = 0$).
Exponential $\mathcal E(\lambda)$: $\varphi(\xi) =
\lambda\int_0^\infty\eu^{(\iu\xi - \lambda)x}\dd x =
\frac{\lambda}{\lambda - \iu\xi}$ (the antiderivative
vanishes at $+\infty$ since $\operatorname{Re}(\iu\xi -
\lambda) < 0$). Poisson $\mathcal P(\lambda)$: by the
transfer theorem for discrete laws,
\[
\varphi(\xi) = \sum_{k\geq0}\eu^{\iu\xi
k}\,\eu^{-\lambda}\frac{\lambda^k}{k!}
= \eu^{-\lambda}\exp\bigl(\lambda\eu^{\iu\xi}\bigr)
= \exp\bigl(\lambda(\eu^{\iu\xi} - 1)\bigr).
\]
Binomial $\mathcal B(n, p)$: a sum of $n$ \omterm{def:b3:probability:independence}{independent}
Bernoulli variables, each with cf $1 - p + p\eu^{\iu\xi}$,
so $\varphi(\xi) = \bigl(1 - p + p\eu^{\iu\xi}\bigr)^n$
(\cref{prop:b3:clt:cfbasics}(b)). Poisson additivity: if $X
\sim \mathcal P(\lambda)$, $Y \sim \mathcal P(\mu)$ are
\omterm{def:b3:probability:independence}{independent},
\[
\varphi_{X+Y}(\xi) = \eu^{\lambda(\eu^{\iu\xi}-1)}
\eu^{\mu(\eu^{\iu\xi}-1)} =
\eu^{(\lambda+\mu)(\eu^{\iu\xi}-1)},
\]
the cf of $\mathcal P(\lambda + \mu)$; injectivity
(\cref{thm:b3:clt:injectivity}) identifies the law.
\end{solution}

\begin{solution}{exo:b3:clt:2}
(a) $\overline{\varphi_X(\xi)} = \E\eu^{-\iu\xi X} =
\varphi_{-X}(\xi)$. So $\varphi_X$ is real if and only if
$\varphi_X = \varphi_{-X}$, if and only if (injectivity,
\cref{thm:b3:clt:injectivity}) $X$ and $-X$ have the same
law.
(b) Write $\varphi_X(\xi_0) = \eu^{\iu\theta}$. Then
\[
\E\bigl[1 - \cos(\xi_0X - \theta)\bigr]
= 1 - \operatorname{Re}\bigl(\eu^{-\iu\theta}
\varphi_X(\xi_0)\bigr) = 1 - 1 = 0 .
\]
The integrand is nonnegative, so $\cos(\xi_0X - \theta) = 1$
almost surely (a nonnegative variable with zero \omterm{def:b3:probability:space}{expectation}
vanishes a.s.), i.e.\ $\xi_0X - \theta \in 2\pi\Z$ a.s.: $X$
takes its values in the arithmetic progression
$\frac{\theta}{\xi_0} + \frac{2\pi}{\xi_0}\Z$ almost surely.
If $X$ has a \omterm{ex:b3:lebesgue:gamma}{density}, this countable set is Lebesgue-null,
so it carries probability $0$ --- contradiction; therefore
$\abs{\varphi_X(\xi)} < 1$ for every $\xi \neq 0$.
\end{solution}

\begin{solution}{exo:b3:clt:3}
By \omterm{def:b3:probability:independence}{independence} and \cref{prop:b3:clt:cfbasics}:
\[
\varphi_{X+Y}(\xi) = \eu^{\iu m_1\xi - \sigma_1^2\xi^2/2}\,
\eu^{\iu m_2\xi - \sigma_2^2\xi^2/2}
= \eu^{\iu(m_1+m_2)\xi - (\sigma_1^2+\sigma_2^2)\xi^2/2},
\]
the cf of $\mathcal N(m_1 + m_2, \sigma_1^2 + \sigma_2^2)$;
injectivity concludes. Stability under affine maps is the
affine rule ($aX + b \sim \mathcal N(am_1 + b,
a^2\sigma_1^2)$, allowing the degenerate case $a = 0$), and
stability under \omterm{def:b3:probability:independence}{independent} sums follows by induction on the
computation above. For \emph{dependent} \omterm{def:b3:clt:gaussianvector}{Gaussians} the sum
need not be \omterm{def:b3:clt:gaussianvector}{Gaussian}: in \cref{exo:b3:clt:9}, $X$ and $Y =
\varepsilon X$ are each standard \omterm{def:b3:clt:gaussianvector}{Gaussian} but $X + Y$
vanishes with probability $\frac12$ without being a.s.\
zero, so it is not \omterm{def:b3:clt:gaussianvector}{Gaussian}.
\end{solution}

\begin{solution}{exo:b3:clt:4}
(a) \emph{Direct implication.} Let $t$ be a \omterm{def:b3:topology:continuity}{continuity} point
of $F_X$ and $\delta > 0$. Take the \omterm{def:b3:topology:continuity}{continuous} ramps
$f^-$ ($= 1$ on $\intoc{-\infty}{t-\delta}$, $0$ from $t$
on, affine between) and $f^+$ ($= 1$ on $\intoc{-\infty}t$,
$0$ from $t + \delta$ on, affine between); then $f^- \leq
\mathbf 1_{\intoc{-\infty}t} \leq f^+$, so
\[
\E f^-(X_n) \leq F_{X_n}(t) \leq \E f^+(X_n),
\]
and the outer terms converge to $\E f^\pm(X)$, themselves
squeezed between $F_X(t - \delta)$ and $F_X(t + \delta)$.
Letting $n \to \infty$ then $\delta \to 0$ and using
\omterm{def:b3:topology:continuity}{continuity} of $F_X$ at $t$: $F_{X_n}(t) \to F_X(t)$.

\emph{Converse.} Let $f$ be bounded \omterm{def:b3:topology:continuity}{continuous}, $M =
\sup\abs f$, $\varepsilon > 0$. The \omterm{def:b3:topology:continuity}{continuity} points of
$F_X$ are dense ($F_X$ has at most countably many jumps), so
choose \omterm{def:b3:topology:continuity}{continuity} points $a < b$ with $F_X(a) < \varepsilon$
and $1 - F_X(b) < \varepsilon$. On the \omterm{def:b3:topology:compact}{compact} $\intcc ab$
the function $f$ is uniformly \omterm{def:b3:topology:continuity}{continuous}: choose \omterm{def:b3:topology:continuity}{continuity}
points $a = t_0 < t_1 < \dots < t_m = b$ of $F_X$ with
oscillation of $f$ at most $\varepsilon$ on each
$\intoc{t_{j-1}}{t_j}$, and set $g = \sum_j
f(t_j)\,\mathbf 1_{\intoc{t_{j-1}}{t_j}}$. Then $\abs{f - g}
\leq \varepsilon$ on $\intoc ab$, $\abs g \leq M$, and for
$T = X_n$ or $X$:
\[
\bigl|\E f(T) - \E g(T)\bigr| \leq \varepsilon +
2M\bigl(F_T(a) + 1 - F_T(b)\bigr).
\]
Moreover $\E g(X_n) = \sum_j f(t_j)\bigl(F_{X_n}(t_j) -
F_{X_n}(t_{j-1})\bigr) \to \E g(X)$ (a finite sum of
converging terms, all $t_j$ being \omterm{def:b3:topology:continuity}{continuity} points), and
$F_{X_n}(a) \to F_X(a) < \varepsilon$, $1 - F_{X_n}(b) \to 1
- F_X(b) < \varepsilon$. Assembling:
$\limsup_n\abs{\E f(X_n) - \E f(X)} \leq 2\varepsilon +
8M\varepsilon$; let $\varepsilon \to 0$.

(b) The distribution function of the constant $c$ is
$\mathbf 1_{\intco c\infty}$, \omterm{def:b3:topology:continuity}{continuous} except at $c$. For
$\varepsilon > 0$, the points $c - \varepsilon$ and $c +
\frac\varepsilon2$ are \omterm{def:b3:topology:continuity}{continuity} points, so
\[
\P(\abs{X_n - c} > \varepsilon) \leq F_{X_n}(c -
\varepsilon) + 1 - F_{X_n}\Bigl(c + \frac\varepsilon2\Bigr)
\longrightarrow 0 + 1 - 1 = 0 .
\]
\end{solution}

\begin{solution}{exo:b3:clt:5}
Let $z_n = p_n(\eu^{\iu\xi} - 1)$, so $\varphi_{X_n}(\xi) =
(1 + z_n)^n$ (\cref{exo:b3:clt:1}) and $\abs{z_n} \leq 2p_n
\to 0$ (note $p_n = \frac{np_n}n \to 0$). Both $1 + z_n$ and
$\eu^{z_n}$ have modulus at most $1$: $\abs{1 + z_n} =
\abs{(1 - p_n) + p_n\eu^{\iu\xi}} \leq 1$ by the triangle
inequality, and $\abs{\eu^{z_n}} = \eu^{p_n(\cos\xi - 1)}
\leq 1$. The telescoping inequality $\abs{a^n - b^n} \leq
n\abs{a - b}$ (proof of \cref{thm:b3:clt:clt}) and the power
series bound $\abs{\eu^z - 1 - z} \leq
\abs z^2\eu^{\abs z}$ give
\[
\bigl|(1 + z_n)^n - \eu^{nz_n}\bigr| \leq n\bigl|1 + z_n -
\eu^{z_n}\bigr| \leq n\,\abs{z_n}^2\,\eu^{\abs{z_n}} \leq
4\eu^2\,np_n^2 = 4\eu^2\,(np_n)\,p_n \longrightarrow 0 .
\]
Since $nz_n = np_n(\eu^{\iu\xi} - 1) \to
\lambda(\eu^{\iu\xi} - 1)$, we conclude $\varphi_{X_n}(\xi)
\to \exp\bigl(\lambda(\eu^{\iu\xi} - 1)\bigr)$ for every
$\xi$: the cf of $\mathcal P(\lambda)$, and L\'evy
(\cref{thm:b3:clt:levy}) gives $X_n \Rightarrow \mathcal
P(\lambda)$. Numerically: $\P\bigl(\mathcal B(100, 0.02) =
0\bigr) = 0.98^{100} = \eu^{100\ln 0.98} \approx
\eu^{-2.020} \approx 0.1326$, while $\P\bigl(\mathcal P(2) =
0\bigr) = \eu^{-2} \approx 0.1353$: two percent apart
already at this coarse $n$.
\end{solution}

\begin{solution}{exo:b3:clt:6}
(a) One roll has mean $\frac72$ and variance
$\frac{35}{12}$, so $S$ has mean $3500$, variance
$\frac{35000}{12} \approx 2916.7$ and standard deviation
$\approx 54.0$. By the CLT,
\[
\P(S > 3600) = \P\Bigl(\frac{S - 3500}{54.0} > 1.85\Bigr)
\approx 1 - \Phi(1.85) \approx 0.032 :
\]
about a $3\%$ chance.
(b) $S \sim \mathcal B(100, \frac12)$: mean $50$, standard
deviation $5$. With the \omterm{def:b3:topology:continuity}{continuity} correction,
\[
\P(45 \leq S \leq 55) \approx
\Phi\Bigl(\frac{55.5 - 50}{5}\Bigr) -
\Phi\Bigl(\frac{44.5 - 50}{5}\Bigr) = 2\Phi(1.1) - 1
\approx 0.729,
\]
against the exact value $0.7287$; without the correction,
$2\Phi(1) - 1 \approx 0.683$, off by almost five points. The
correction matters because $S$ is a lattice variable: the
atom $\P(S = k)$ is well approximated by the \omterm{def:b3:clt:gaussianvector}{Gaussian} mass
of $\intcc{k - \frac12}{k + \frac12}$, and clipping the
interval at the integers $45$ and $55$ discards half an atom
at each end.
\end{solution}

\begin{solution}{exo:b3:clt:7}
The variables $g(U_k)$ are i.i.d.\ (\omterm{def:b3:lebesgue:measurable}{measurable} images of
i.i.d.\ variables), \omterm{def:b3:lebesgue:l1}{square-integrable}, with mean $I$
(transfer theorem, \cref{exo:b3:product:9}) and variance
$\sigma^2$. If $\sigma > 0$, \cref{thm:b3:clt:clt} applied
to them is exactly the stated convergence
\[
\sqrt n\,\Bigl(\frac1n\sum_{k\leq n}g(U_k) - I\Bigr) =
\frac{\sum_{k\leq n}\bigl(g(U_k) - I\bigr)}{\sqrt n}
\Longrightarrow \mathcal N(0, \sigma^2)
\]
(if $\sigma = 0$, $g$ is a.s.\ constant and the left-hand
side vanishes identically). Hence
$\P\bigl(\abs{\frac1n\sum g(U_k) - I} \leq
1.96\,\sigma/\sqrt n\bigr) \to 0.95$: the error bar $\pm
1.96\,\sigma/\sqrt n$ sees the dimension $d$ only through
the constant $\sigma$, never through the rate in $n$. The
midpoint rule with $n$ nodes in dimension $d$ has mesh
$n^{-1/d}$ and error of order $n^{-2/d}$ for $\mathcal C^2$
integrands. \omterm{ex:b3:probability:sllnapps}{Monte Carlo}'s $n^{-1/2}$ decays faster than
$n^{-2/d}$ exactly when $\frac12 > \frac2d$, i.e.\ $d > 4$:
from dimension $5$ on, random sampling asymptotically beats
the grid --- the curse of dimensionality spares
probabilistic methods, which is why \omterm{ex:b3:probability:sllnapps}{Monte Carlo} rules
high-dimensional integration.
\end{solution}

\begin{solution}{exo:b3:clt:8}
\emph{Sum.} For fixed $\xi$:
\[
\bigl|\E\eu^{\iu\xi(X_n+Y_n)} -
\eu^{\iu\xi c}\,\E\eu^{\iu\xi X_n}\bigr|
= \bigl|\E\bigl[\eu^{\iu\xi X_n}\bigl(\eu^{\iu\xi Y_n} -
\eu^{\iu\xi c}\bigr)\bigr]\bigr|
\leq \E\bigl|\eu^{\iu\xi(Y_n - c)} - 1\bigr| .
\]
Split on the event $\{\abs{Y_n - c} \leq \delta\}$: there,
$\abs{\eu^{\iu\xi(Y_n-c)} - 1} \leq \abs\xi\,\delta$ (the
chord is shorter than the arc); the complement contributes
at most $2\,\P(\abs{Y_n - c} > \delta) \to 0$. Hence the
$\limsup$ is $\leq \abs\xi\,\delta$ for every $\delta > 0$:
the difference tends to $0$. Since $\E\eu^{\iu\xi X_n} \to
\varphi_X(\xi)$, we get $\varphi_{X_n+Y_n}(\xi) \to
\eu^{\iu\xi c}\varphi_X(\xi) = \varphi_{X+c}(\xi)$, and
L\'evy (\cref{thm:b3:clt:levy}) yields $X_n + Y_n
\Rightarrow X + c$.

\emph{Product.} First, $cX_n \Rightarrow cX$:
$\varphi_{cX_n}(\xi) = \varphi_{X_n}(c\xi) \to
\varphi_X(c\xi) = \varphi_{cX}(\xi)$. Next, $(Y_n - c)X_n
\to 0$ in probability: the laws of the $X_n$ are tight (their
cfs converge to a cf; see the tightness step of
\cref{thm:b3:clt:levy}), so given $\varepsilon > 0$ pick $M$
with $\P(\abs{X_n} > M) \leq \varepsilon$ for all $n$; then
\[
\P\bigl(\abs{(Y_n - c)X_n} > \varepsilon\bigr) \leq
\P(\abs{X_n} > M) + \P\Bigl(\abs{Y_n - c} >
\frac{\varepsilon}{M}\Bigr) \leq \varepsilon + o(1) .
\]
Writing $Y_nX_n = cX_n + (Y_n - c)X_n$ and applying the sum
part (whose proof only used $Y_n' := (Y_n - c)X_n \to 0$ in
probability, with constant $0$): $Y_nX_n \Rightarrow cX$.

\emph{Application.} By the strong law of large numbers
(\cref{thm:b3:probability:slln}), $\hat p_n \to p$ a.s., so
by \omterm{def:b3:topology:continuity}{continuity} $\hat\sigma_n = \sqrt{\hat p_n(1 - \hat p_n)}
\to \sigma = \sqrt{p(1 - p)} > 0$ a.s., hence
$\frac{\sigma}{\hat\sigma_n} \to 1$ in probability. Slutsky's
product rule upgrades $\frac{S_n - np}{\sigma\sqrt n}
\Rightarrow \mathcal N(0,1)$ to $\frac{S_n -
np}{\hat\sigma_n\sqrt n} = \frac{\sigma}{\hat\sigma_n}\cdot
\frac{S_n - np}{\sigma\sqrt n} \Rightarrow \mathcal N(0,1)$:
the \emph{usable} confidence interval $\hat p_n \pm
1.96\,\hat\sigma_n/\sqrt n$, built from the data alone,
keeps its asymptotic $95\%$ level.
\end{solution}

\begin{solution}{exo:b3:clt:9}
(a) Splitting the \omterm{def:b3:probability:space}{expectation} over the two values of
$\varepsilon$ (\omterm{def:b3:probability:independence}{independence}): for Borel $B$, $\P(Y \in B) =
\frac12\P(X \in B) + \frac12\P(-X \in B) = \P(X \in B)$,
since $-X \sim X$ ($\mathcal N(0,1)$ is symmetric): $Y \sim
\mathcal N(0,1)$. And $\operatorname{Cov}(X, Y) =
\E[\varepsilon X^2] = \E[\varepsilon]\,\E[X^2] = 0 \cdot 1 =
0$.
(b) $\abs Y = \abs X$, so $\P(\abs X \leq 1,\ \abs Y \geq 2)
= 0$ while $\P(\abs X \leq 1)\,\P(\abs Y \geq 2) > 0$: not
\omterm{def:b3:probability:independence}{independent}. If $(X, Y)$ were a \omterm{def:b3:clt:gaussianvector}{Gaussian vector}, $X + Y =
(1 + \varepsilon)X$ would be a real \omterm{def:b3:clt:gaussianvector}{Gaussian} variable
(\cref{def:b3:clt:gaussianvector} with $t = (1,1)$); but
$\P(X + Y = 0) = \P(\varepsilon = -1) = \frac12$, whereas a
\omterm{def:b3:clt:gaussianvector}{Gaussian} variable has an atom only if it is a.s.\ constant
--- and $X + Y$ equals $2X \neq 0$ a.s.\ on
$\{\varepsilon = 1\}$. Contradiction: $(X, Y)$ is not
\omterm{def:b3:clt:gaussianvector}{Gaussian}.
(c) Each marginal is \omterm{def:b3:clt:gaussianvector}{Gaussian} and the covariance vanishes,
yet \omterm{def:b3:probability:independence}{independence} fails --- because the \emph{pair} is not
jointly \omterm{def:b3:clt:gaussianvector}{Gaussian}. \cref{thm:b3:clt:gaussianvector}(2) cannot
be weakened to ``\omterm{def:b3:clt:gaussianvector}{Gaussian} marginals''.
\end{solution}

\begin{solution}{exo:b3:clt:10}
(a) \cref{exo:b3:fouriertransform:1} computes
$\widehat{\eu^{-\abs\cdot}}(\xi) = \frac{2}{1 + \xi^2}$;
both sides being \omterm{def:b3:lebesgue:l1}{integrable}, Fourier inversion
(\cref{thm:b3:fouriertransform:inversion}) turns this
around:
\[
\int_\R\eu^{\iu\xi x}\,\frac{\dd x}{\pi(1 + x^2)} =
\eu^{-\abs\xi},
\]
which is exactly $\varphi_X(\xi)$ for a Cauchy variable $X$.
(b) By \omterm{def:b3:probability:independence}{independence}, $\varphi_{S_n}(\xi) =
\bigl(\eu^{-\abs\xi}\bigr)^n = \eu^{-n\abs\xi}$, so
$\varphi_{S_n/n}(\xi) = \varphi_{S_n}(\xi/n) =
\eu^{-\abs\xi}$: the empirical mean $\frac{S_n}n$ is again
standard Cauchy for every $n$ (injectivity). The average
never concentrates: its fluctuations at time $10^6$ are
those of a single observation.
(c) $\E\abs{X_1} = \frac2\pi\int_0^\infty\frac{x\,\dd x}{1 +
x^2} = +\infty$: the Cauchy law is not \omterm{def:b3:lebesgue:l1}{integrable}, so the
strong law of large numbers
(\cref{thm:b3:probability:slln}) does not apply, and the CLT
(which needs a finite variance) even less. Here their
conclusions genuinely fail, not merely their proofs.
Consistency check: $\varphi(\xi) = \eu^{-\abs\xi}$ is not
differentiable at $0$, as \cref{prop:b3:clt:cfbasics}(c)
read contrapositively predicts for a \omterm{def:b3:lebesgue:l1}{non-integrable}
variable.
\end{solution}

\begin{solution}{exo:b3:clt:11}
(a) $\varphi_{aX_1 + bX_2}(\xi) =
\eu^{-a\abs\xi}\eu^{-b\abs\xi} = \eu^{-(a+b)\abs\xi} =
\varphi_{(a+b)X_1}(\xi)$ (\omterm{def:b3:probability:independence}{independence} and
\cref{exo:b3:clt:10}); injectivity identifies the laws.

(b) $\varphi_{aX_1+bX_2}(\xi) = \eu^{-a^2\xi^2/2}
\eu^{-b^2\xi^2/2} = \eu^{-(a^2+b^2)\xi^2/2}$: the law of
$\sqrt{a^2+b^2}\,X_1$.

(c) If $\varphi_X(\xi) = \eu^{-c\abs\xi^\alpha}$, then
$S_n = X_1 + \dots + X_n$ has $\varphi_{S_n} =
\eu^{-cn\abs\xi^\alpha}$, and $S_n/n^{1/\alpha}$ has
$\varphi(\xi) = \eu^{-c\abs\xi^\alpha}$ again: exact
self-reproduction under the $n^{1/\alpha}$ scaling ---
$\sqrt n$ for the \omterm{def:b3:clt:gaussianvector}{Gaussian} ($\alpha = 2$), $n$ itself for
Cauchy ($\alpha = 1$, \cref{exo:b3:clt:10}(b)). A sum of
i.i.d.\ variables can only converge (after affine
normalization) to a law that is stable under such
convolutions; the CLT says finite variance forces the
\omterm{def:b3:clt:gaussianvector}{Gaussian} basin, and the Cauchy basin is reserved for laws
with tails so heavy that $\E X^2 = \infty$ and even $\E\abs
X = \infty$ --- e.g.\ sums of Cauchy variables themselves.
Universality has several islands, indexed by the tail
exponent $\alpha \in \intoc02$.
\end{solution}

\begin{solution}{exo:b3:clt:12}
(a) The indicators $\mathbf 1_{X_k \leq t}$ are i.i.d.\
Bernoulli of parameter $p = F(t)$: their sum $nF_n(t)$ is
binomial $\mathcal B(n, p)$; the strong law gives $F_n(t)
\to p$ a.s., and the CLT (\cref{thm:b3:clt:clt}) applied to
the same indicators (variance $p(1-p)$) gives the stated
\omterm{def:b3:clt:gaussianvector}{Gaussian} limit.

(b) $p(1 - p)$ is maximal at $p = \frac12$, i.e.\ where
$F(t) = \frac12$: at the \emph{median}. Estimating tail
probabilities is asymptotically easy (variance $\to 0$ as
$p \to 0, 1$); the median region carries the largest
statistical noise --- the empirical curve wobbles most in
its middle.

(c) For \omterm{def:b3:topology:continuity}{continuous} $F$, the variables $U_k = F(X_k)$ are
i.i.d.\ uniform on $\intoo01$ (\cref{exo:b3:probability:1}),
and monotonicity of $F$ gives, writing $G_n$ for the
empirical distribution function of the $U_k$:
\[
\sup_{t\in\R}\,\abs{F_n(t) - F(t)}
= \sup_{u \in \operatorname{im}F}\,\abs{G_n(u) - u}
= \sup_{u\in\intcc01}\abs{G_n(u) - u} :
\]
the first equality because $\{X_k \leq t\} = \{U_k \leq
F(t)\}$ up to null events (monotonicity; strict inequality
can fail only on the flat parts of $F$, where both sides
are unchanged), and the second because a \omterm{def:b3:topology:continuity}{continuous} $F$,
running from $0$ to $1$, attains every value of
$\intoo01$ (intermediate value theorem), and the endpoints
add nothing ($G_n(0) - 0 = 0$ and $G_n(1) - 1 = 0$). The
right side involves only uniforms: one law
for all $F$ --- so a single table of critical values (that
of Kolmogorov's distribution) tests \emph{any} \omterm{def:b3:topology:continuity}{continuous}
model against data.
\end{solution}

\begin{solution}{pb:b3:clt:1}
\textbf{1.} The family $(X_1, \dots, X_n, N_1, \dots, N_n)$
is \omterm{def:b3:probability:independence}{independent}: the two blocks are \omterm{def:b3:probability:independence}{independent} of each
other by construction and each block is i.i.d. $W_i$ is a
\omterm{def:b3:lebesgue:measurable}{measurable function} of the variables $(X_j)_{j<i}$ and
$(N_j)_{j>i}$ only, all distinct from $X_i$ and $N_i$: by
the coalition principle
(\cref{thm:b3:probability:independence}), $W_i$ is
\omterm{def:b3:probability:independence}{independent} of the pair $(X_i, N_i)$. The decompositions
$H_i = W_i + \frac{X_i}{\sqrt n}$ and $H_{i-1} = W_i +
\frac{N_i}{\sqrt n}$ are immediate from the definitions:
passing from $H_i$ to $H_{i-1}$ swaps the single summand
$X_i$ for $N_i$.

\textbf{2.} Taylor--Lagrange at order $3$: there is $c$
between $w$ and $w + h$ with $f(w + h) = f(w) + f'(w)h +
\frac12f''(w)h^2 + \frac16f'''(c)h^3$, and $\abs{f'''(c)}
\leq M_3$ gives the bound.

\textbf{3.} Subtracting the two expansions at the common
base point $w = W_i$:
\[
f(H_i) - f(H_{i-1}) = f'(W_i)\,\frac{X_i - N_i}{\sqrt n} +
\frac{f''(W_i)}{2}\,\frac{X_i^2 - N_i^2}{n} + R_i,
\qquad
\abs{R_i} \leq \frac{M_3}{6}\cdot
\frac{\abs{X_i}^3 + \abs{N_i}^3}{n^{3/2}} .
\]
Take \omterm{def:b3:probability:space}{expectations}. By question 1, $f'(W_i)$ and $f''(W_i)$
are \omterm{def:b3:probability:independence}{independent} of $(X_i, N_i)$, so the mixed \omterm{def:b3:probability:space}{expectations}
factor:
\begin{align*}
\E\Bigl[f'(W_i)\,\frac{X_i - N_i}{\sqrt n}\Bigr] &=
\E\bigl[f'(W_i)\bigr]\,\frac{\E X_i - \E N_i}{\sqrt n} = 0,
\\
\E\Bigl[f''(W_i)\,\frac{X_i^2 - N_i^2}{n}\Bigr] &=
\E\bigl[f''(W_i)\bigr]\,\frac{1 - 1}{n} = 0 :
\end{align*}
the first two moments of $X_i$ and $N_i$ \emph{match}, and
only the remainder survives:
\[
\bigl|\E f(H_i) - \E f(H_{i-1})\bigr| \leq \E\abs{R_i} \leq
\frac{M_3}{6}\cdot\frac{\beta + \gamma}{n^{3/2}} .
\]
The \omterm{def:b3:clt:gaussianvector}{Gaussian} third moment: $\gamma = \E\abs{N_1}^3 =
2\int_0^\infty x^3\,\frac{\eu^{-x^2/2}}{\sqrt{2\pi}}\,\dd x
= \frac{2}{\sqrt{2\pi}}\int_0^\infty 2u\,\eu^{-u}\dd u =
\frac{4}{\sqrt{2\pi}} = \frac{2\sqrt2}{\sqrt\pi}$
(substitution $u = x^2/2$, then $\Gamma(2) = 1$).

\textbf{4.} Telescoping $\E f(T_n) - \E f(G_n) =
\sum_{i=1}^n\bigl(\E f(H_i) - \E f(H_{i-1})\bigr)$ and
applying question 3 to each of the $n$ terms:
\[
\bigl|\E f(T_n) - \E f(G_n)\bigr| \leq
n \cdot \frac{M_3(\beta + \gamma)}{6\,n^{3/2}} =
\frac{M_3\,(\beta + \gamma)}{6\,\sqrt n} .
\]

\textbf{5.} $G_n$ is $\mathcal N(0,1)$ \emph{exactly} for
every $n$ (a normalized sum of \omterm{def:b3:probability:independence}{independent} standard
\omterm{def:b3:clt:gaussianvector}{Gaussians}, \cref{exo:b3:clt:3}), so $\E f(G_n) = \E f(N)$
and question 4 reads $\abs{\E f(T_n) - \E f(N)} \leq
\frac{M_3(\beta+\gamma)}{6\sqrt n} \to 0$ for $f \in
\mathcal C^3_b$. \emph{Upgrade.} Let $f$ be bounded
\omterm{def:b3:topology:continuity}{continuous}, $M = \sup\abs f$, $\varepsilon > 0$. Choose $A
\geq 1$ with $\frac1{A^2} \leq \varepsilon$: Chebyshev with
$\V(T_n) = 1$ gives $\P(\abs{T_n} > A) \leq \varepsilon$
for all $n$, and likewise $\P(\abs N > A) \leq \varepsilon$.
Let $\chi$ be $\mathcal C^\infty$ with
$\mathbf 1_{\intcc{-A}A} \leq \chi \leq
\mathbf 1_{\intcc{-A-1}{A+1}}$ (a smooth plateau, built by
mollifying $\mathbf 1_{\intcc{-A-\frac12}{A+\frac12}}$,
\cref{thm:b3:lp:regularization}); $g = f\chi$ is \omterm{def:b3:topology:continuity}{continuous}
with \omterm{def:b3:topology:compact}{compact} support, hence uniformly \omterm{def:b3:topology:continuity}{continuous}, so its
mollification $g_\eta = g * \rho_\eta$ is $\mathcal
C^\infty$ with bounded derivatives of all orders and
$\norm{g - g_\eta}_\infty \leq \varepsilon$ for $\eta$
small enough. For $T = T_n$ or $N$, since $f = g$ on
$\intcc{-A}A$ and $\abs{f - g} \leq 2M$ everywhere:
\[
\bigl|\E f(T) - \E g_\eta(T)\bigr| \leq
\E\abs{(f - g)(T)} + \norm{g - g_\eta}_\infty
\leq 2M\,\P(\abs T > A) + \varepsilon \leq
(2M + 1)\,\varepsilon .
\]
Combining with $\E g_\eta(T_n) \to \E g_\eta(N)$ (question
4 applies: $g_\eta \in \mathcal C^3_b$):
\[
\limsup_n\;\bigl|\E f(T_n) - \E f(N)\bigr| \leq
2(2M + 1)\,\varepsilon ,
\]
and $\varepsilon$ was arbitrary: $\E f(T_n) \to \E f(N)$ for
every bounded \omterm{def:b3:topology:continuity}{continuous} $f$, i.e.\ $T_n \Rightarrow
\mathcal N(0,1)$.

\textbf{6.} Identical distribution entered only through one
sentence: ``$X_i$ and $N_i$ have the same first two
moments''. So let $X_1, \dots, X_n$ be \omterm{def:b3:probability:independence}{independent},
centered, with variances $\sigma_i^2$ and finite third
moments, $s_n^2 = \sum_i\sigma_i^2 > 0$, and take $N_i \sim
\mathcal N(0, \sigma_i^2)$ \omterm{def:b3:probability:independence}{independent} of everything.
Define the hybrids with normalization $s_n$: $H_i =
\frac1{s_n}(\sum_{j\leq i}X_j + \sum_{j>i}N_j)$. In the
$i$-th swap, $\E X_i = \E N_i = 0$ and $\E X_i^2 = \E N_i^2
= \sigma_i^2$ again kill the $f'$ and $f''$ terms, and the
remainder gives (using $\E\abs{N_i}^3 = \sigma_i^3\gamma$ by
scaling):
\[
\bigl|\E f(H_i) - \E f(H_{i-1})\bigr| \leq
\frac{M_3}{6\,s_n^3}\bigl(\E\abs{X_i}^3 +
\sigma_i^3\gamma\bigr) .
\]
Telescoping:
\[
\Bigl|\E f\Bigl(\frac{X_1 + \dots + X_n}{s_n}\Bigr) -
\E f(N)\Bigr| \leq \frac{M_3}{6\,s_n^3}\sum_{i=1}^n
\Bigl(\E\abs{X_i}^3 + \V(X_i)^{3/2}\,\gamma\Bigr) .
\]
Since $\sigma_i^3 = (\E X_i^2)^{3/2} \leq \E\abs{X_i}^3$
(the power--mean inequality, i.e.\ Jensen for $t \mapsto
t^{3/2}$ applied to $X_i^2$), the right side is at most
$\frac{M_3(1 + \gamma)}{6}\cdot
\frac{\sum_i\E\abs{X_i}^3}{s_n^3}$: under \emph{Lyapunov's
condition} $\frac1{s_n^3}\sum_i\E\abs{X_i}^3 \to 0$, the
normalized sums converge in law to $\mathcal N(0,1)$ ---
the CLT without identical distribution.

\textbf{7.} Let $\rho \in \mathcal C^\infty_c(\intoo01)$
with $\int\rho = 1$ and set $\psi(x) = \int_x^1\rho(s)\dd
s$: $\psi$ is $\mathcal C^\infty$, nonincreasing, $\psi = 1$
on $\R_-$, $\psi = 0$ on $\intco1\infty$; let $K =
\norm{\psi'''}_\infty$. For $t \in \R$ and $\delta > 0$
define $\psi_\delta(x) = \psi\bigl(\frac{x - t}\delta\bigr)$
and $\tilde\psi_\delta(x) = \psi\bigl(\frac{x - t}\delta +
1\bigr)$: these are $\mathcal C^3_b$ with third derivative
bounded by $K/\delta^3$, and
\[
\mathbf 1_{\intoc{-\infty}{t-\delta}} \leq
\tilde\psi_\delta \leq \mathbf 1_{\intoc{-\infty}t} \leq
\psi_\delta \leq \mathbf 1_{\intoc{-\infty}{t+\delta}} .
\]
Upper bound: by question 4 applied to $\psi_\delta$ (with
$M_3 = K/\delta^3$),
\[
\P(T_n \leq t) \leq \E\psi_\delta(T_n) \leq
\E\psi_\delta(N) + \frac{K(\beta +
\gamma)}{6\,\delta^3\sqrt n}
\leq \Phi(t + \delta) + \frac{K(\beta +
\gamma)}{6\,\delta^3\sqrt n}
\leq \Phi(t) + \frac{\delta}{\sqrt{2\pi}} +
\frac{K(\beta + \gamma)}{6\,\delta^3\sqrt n},
\]
because $\Phi$ is Lipschitz with constant
$\frac1{\sqrt{2\pi}}$ (its \omterm{ex:b3:lebesgue:gamma}{density} is bounded by
$\frac1{\sqrt{2\pi}}$). The symmetric lower bound via
$\tilde\psi_\delta$ gives the two-term estimate
\[
\sup_{t\in\R}\,\bigl|\P(T_n \leq t) - \Phi(t)\bigr| \leq
\frac{K(\beta + \gamma)}{6}\cdot\frac{1}{\delta^3\sqrt n} +
\frac{\delta}{\sqrt{2\pi}}
\qquad(\delta > 0\ \text{arbitrary}).
\]
The two terms balance when $\delta^{-3}n^{-1/2} \asymp
\delta$, i.e.\ $\delta = n^{-1/8}$: both are then
$O(n^{-1/8})$, an explicit uniform rate valid for every
$n$. (The optimal Berry--Esseen rate $C\beta/\sqrt n$
requires the Fourier smoothing inequality; swapping trades
sharpness for \omterm{def:b3:complete:complete}{complete} elementarity.)

\textbf{8.} For $X_i = 2B_i - 1$ (fair signs): centered,
variance $1$, and $\abs{X_i} = 1$ so $\beta = 1$. Question 7
then bounds $\sup_t\abs{\P(\frac{S_n}{\sqrt n} \leq t) -
\Phi(t)}$ explicitly and uniformly for \emph{every} finite
$n$ --- a global, non-asymptotic statement about the
distribution function. The local estimate of
\cref{pb:b3:product:1}, question 7, gives instead the exact
asymptotics of an individual atom, $\P(S_{2n} = 2k) \sim
\frac{\eu^{-k^2/n}}{\sqrt{\pi n}}$: it resolves
probabilities of size $n^{-1/2}$, far below question 7's
$n^{-1/8}$ resolution, but it is pointwise, asymptotic
(no explicit error at fixed $n$) and tied to this
particular lattice law. Local precision versus global
uniformity: the two methods are complementary, and summing
the local estimate over $k \in \intint{a\sqrt n}{b\sqrt n}$
recovers de Moivre--Laplace on intervals --- with a sharper
rate, but only for this law.

\textbf{9.} The swapping argument used nothing about the law
of the $X_i$ beyond $\E X_i = 0$, $\E X_i^2 = 1$ and the
finiteness of $\E\abs{X_i}^3$: had we replaced the \omterm{def:b3:clt:gaussianvector}{Gaussians}
$N_i$ by any other i.i.d.\ family with the same first two
moments and finite third moment, the same telescoping would
bound $\abs{\E f(\text{sum}_X) - \E f(\text{sum}_Y)}$ by
$O(n^{-1/2})$ for every smooth $f$. Smooth statistics of
large \omterm{def:b3:probability:independence}{independent} sums are therefore \emph{universal}: up to
a quantified error, they depend on the summands' law only
through two numbers. This is the invariance principle:
prove a limit theorem for the most computable law (the
\omterm{def:b3:clt:gaussianvector}{Gaussian}, where everything is exact), then transfer it to
all laws by swapping. The same scheme --- with sums replaced
by more elaborate functionals --- drives Wigner's semicircle
law for random matrices, the universality of roots of random
polynomials, and much of modern probability; the central
limit theorem is its first and simplest instance.

\textbf{10.} $Z$ is a Borel function of $(N_{i+1}, \dots,
N_n)$ only, while $A = W_i + \theta h - Z$ and $h$ are
functions of the remaining variables of the \omterm{def:b3:probability:independence}{independent}
family $(X_1, \dots, X_n, N_1, \dots, N_n)$: by the
coalition principle (\cref{thm:b3:probability:independence}),
$Z$ is \omterm{def:b3:probability:independence}{independent} of $(A, h)$. As a sum of the \omterm{def:b3:probability:independence}{independent}
$N_j/\sqrt n \sim \mathcal N(0, \frac1n)$, $Z \sim \mathcal
N(0, s^2)$ with $s^2 = \frac{n-i}n$ (\cref{exo:b3:clt:3}),
with \omterm{ex:b3:lebesgue:gamma}{density} bounded by $\frac1{s\sqrt{2\pi}}$. The law of
$((A, h), Z)$ is the product of the two marginal laws, so
Tonelli (transfer) freezes the first block: with $G(a) =
\E\abs{g(a + Z)} = \int\abs{g(a + z)}\,\varphi_s(z)\,\dd z
\leq \frac{\norm g_{L^1}}{s\sqrt{2\pi}}$ for every $a$,
\[
\E\bigl[\abs h^3\abs{g(A + Z)}\bigr] =
\E\bigl[\abs h^3\,G(A)\bigr] \leq
\frac{\norm g_{L^1}}{s\sqrt{2\pi}}\,\E\abs h^3
= \sqrt{\frac{n}{2\pi(n-i)}}\,\norm g_{L^1}\,\E\abs h^3 .
\]

\textbf{11.} The integral form of Taylor's formula follows
by integrating $f(w + h) - f(w) =
h\int_0^1f'(w + \theta h)\,\dd\theta$ by parts twice in
$\theta$. Taking \omterm{def:b3:probability:space}{expectations} in the $i$-th swap, the
orders $0, 1, 2$ cancel exactly as in question 3, and the
two remainders (for $h = X_i/\sqrt n$
and $N_i/\sqrt n$) are bounded, for $i \leq n - 1$, by
question 10 with $g = f'''$:
\[
\bigl|\E f(H_i) - \E f(H_{i-1})\bigr| \leq
\int_0^1\frac{(1-\theta)^2}2\,\dd\theta\;
\sqrt{\frac{n}{2\pi(n-i)}}\,\norm{f'''}_{L^1}
\frac{\beta + \gamma}{n^{3/2}}
= \frac{\beta + \gamma}{6\,n^{3/2}}
\sqrt{\frac{n}{2\pi(n-i)}}\,\norm{f'''}_{L^1} .
\]
Summing, with $\sum_{i=1}^{n-1}\sqrt{\frac n{n-i}} =
\sqrt n\sum_{m=1}^{n-1}m^{-1/2} \leq 2n$, and adding
question 3's bound for the last swap ($i = n$, no \omterm{def:b3:clt:gaussianvector}{Gaussian}
left):
\[
\bigl|\E f(T_n) - \E f(G_n)\bigr| \leq
\frac{\beta + \gamma}{3\sqrt{2\pi}}\cdot
\frac{\norm{f'''}_{L^1}}{\sqrt n} +
\frac{M_3(\beta + \gamma)}{6\,n^{3/2}} .
\]
Ramps: $\psi_\delta'''(x) =
\delta^{-3}\psi'''\bigl(\frac{x - t}\delta\bigr)$, so
$M_3 = K\delta^{-3}$ with $K = \norm{\psi'''}_\infty$ and
$\norm{\psi_\delta'''}_{L^1} = \delta^{-2}
\norm{\psi'''}_{L^1} = K_1\delta^{-2}$ (substitution). The
sandwich of question 7 then gives
\[
\sup_t\,\bigl|\P(T_n \leq t) - \Phi(t)\bigr| \leq
\frac{K_1(\beta + \gamma)}{3\sqrt{2\pi}}\cdot
\frac1{\delta^2\sqrt n} +
\frac{K(\beta + \gamma)}{6}\cdot
\frac1{\delta^3n^{3/2}} + \frac\delta{\sqrt{2\pi}} .
\]
At $\delta = n^{-1/6}$ the first and third terms are
$O(n^{-1/6})$ and the middle one $O(n^{-1})$: a uniform
rate $O(n^{-1/6})$, strictly better than question 7's
$n^{-1/8}$ --- the \omterm{def:b3:clt:gaussianvector}{Gaussian} half of the hybrid did the
extra smoothing.

\textbf{12.} $\E N_1^3 = 0$ (odd integrand), and
integration by parts gives $\E N_1^4 = 3\,\E N_1^2 = 3$
($\int x^3\cdot x\varphi(x)\dd x =
3\int x^2\varphi$). For $f$ of class $\mathcal C^4$ with
bounded derivatives, expand each swap to fourth order: the
third-order terms carry the factor $\E X_i^3 - \E N_i^3 =
0$ (\omterm{def:b3:probability:independence}{independence} factors them as in question 3), so only
the fourth-order remainder
$\int_0^1\frac{(1-\theta)^3}6f^{(4)}(w + \theta
h)h^4\dd\theta$ survives, with $\int_0^1
\frac{(1-\theta)^3}6\dd\theta = \frac1{24}$ and $\E h^4 =
\beta_4n^{-2}$ or $3n^{-2}$. Question 10 (with $g =
f^{(4)}$) bounds the swaps $i \leq n - 1$, and summing as
in question 11:
\[
\bigl|\E f(T_n) - \E f(G_n)\bigr| \leq
\frac{\beta_4 + 3}{12\sqrt{2\pi}}\cdot
\frac{\norm{f^{(4)}}_{L^1}}{n} +
\frac{M_4(\beta_4 + 3)}{24\,n^2} .
\]
With $\norm{\psi_\delta^{(4)}}_{L^1} = K_2\delta^{-3}$ and
$M_4 = K'\delta^{-4}$, the distribution-function bound
becomes $C\delta^{-3}n^{-1} + C'\delta^{-4}n^{-2} +
\frac\delta{\sqrt{2\pi}}$; at $\delta = n^{-1/4}$ the
outer terms are $O(n^{-1/4})$ and the middle $O(n^{-1})$:
rate $O(n^{-1/4})$.

\textbf{13.} With $k$ matched moments the surviving
remainder per swap is of order $\E\abs h^{k+1} \asymp
n^{-(k+1)/2}$; the \omterm{def:b3:clt:gaussianvector}{hidden-Gaussian} bound charges
$\norm{f^{(k+1)}}_{L^1}$ and the sum over swaps contributes
the factor $2n$, giving $\asymp\norm{f^{(k+1)}}_{L^1}\,
n^{-(k-1)/2}$ for smooth $f$. Ramps cost
$\norm{\psi_\delta^{(k+1)}}_{L^1} \asymp \delta^{-k}$, so
the distribution-function error is $\asymp
\delta^{-k}n^{-(k-1)/2} + \delta$, balanced at $\delta =
n^{-(k-1)/(2k+2)}$: rate $n^{-(k-1)/(2k+2)}$, which is
$n^{-1/6}$ for $k = 2$, $n^{-1/4}$ for $k = 3$, and tends
to $n^{-1/2}$ only as $k \to \infty$ --- but $k \geq 4$
would force $\E X_1^4 = 3$ and beyond, i.e.\ a law that
already imitates the \omterm{def:b3:clt:gaussianvector}{Gaussian}. The saturation is
structural: swapping adds $n$ swap errors \emph{in
absolute value}, renouncing all cancellation between
swaps. The Fourier proof compares \omterm{def:b3:clt:cf}{characteristic functions},
where the errors appear with their oscillating phases;
Esseen's smoothing inequality converts $\abs{\varphi_{T_n}
- \varphi_N}$, integrated against $\frac{\dd\xi}{\abs\xi}$,
into a distribution-function bound at only logarithmic
cost, and delivers Berry--Esseen's $C\beta n^{-1/2}$ from
three moments. Replacement trades optimality for
robustness --- and, as Part VII shows, for portability.

\textbf{14.} $\Sigma$ is symmetric positive semidefinite;
with $\Sigma = PDP^{\mathsf T}$ ($P$ orthogonal, $D \geq 0$
diagonal, \cref{exo:b3:submanifolds:8}), the symmetric $C =
P\sqrt DP^{\mathsf T}$ satisfies $C^2 = \Sigma$. For any $t
\in \R^2$, $\langle t, CZ\rangle = \langle Ct, Z\rangle =
(Ct)_1Z^1 + (Ct)_2Z^2$ is a linear combination of
\omterm{def:b3:probability:independence}{independent} \omterm{def:b3:clt:gaussianvector}{Gaussians}, hence \omterm{def:b3:clt:gaussianvector}{Gaussian}
(\cref{exo:b3:clt:3}): $N = CZ$ is a \omterm{def:b3:clt:gaussianvector}{Gaussian vector}; its
mean is $0$ and its covariance $\E[NN^{\mathsf T}] =
C\,\E[ZZ^{\mathsf T}]\,C^{\mathsf T} = CC^{\mathsf T} =
\Sigma$. Moments: $\norm N^3 \leq (\abs{N_1} +
\abs{N_2})^3 \leq 4(\abs{N_1}^3 + \abs{N_2}^3)$ (convexity
of $x^3$ on $\R_+$), and each coordinate is a real
\omterm{def:b3:clt:gaussianvector}{Gaussian} with moments of all orders
(\cref{exo:b3:product:10}): $\gamma' < \infty$. Finally
each $\langle t, G_n\rangle = \frac1{\sqrt
n}\sum_i\langle t, N_i\rangle$ is a normalized sum of
i.i.d.\ $\mathcal N(0, t^{\mathsf T}\Sigma t)$, hence
exactly $\mathcal N(0, t^{\mathsf T}\Sigma t)$: $G_n$ is a
\omterm{def:b3:clt:gaussianvector}{Gaussian vector} with mean $0$ and covariance $\Sigma$, and
its law is $\mathcal N(0, \Sigma)$
(\cref{def:b3:clt:gaussianvector}: the law is determined
by these data).

\textbf{15.} Let $\phi(t) = f(w + th)$, $t \in \intcc01$:
$\phi$ is $\mathcal C^3$ with
\[
\phi'''(t) = \sum_{j,k,l\in\{1,2\}}\partial_{jkl}f(w +
th)\,h_jh_kh_l, \qquad \abs{\phi'''(t)} \leq
M_3\Bigl(\sum_j\abs{h_j}\Bigr)^3 = M_3(\abs{h_1} +
\abs{h_2})^3 .
\]
Taylor--Lagrange at order $3$ for $\phi$ between $0$ and
$1$ gives the first inequality; Cauchy--Schwarz gives
$\abs{h_1} + \abs{h_2} \leq \sqrt2\norm h$, whence the
constant $\frac{2\sqrt2M_3}6 = \frac{\sqrt2M_3}3$.

\textbf{16.} Define $H_i$ and $W_i$ as in question 1, now
in $\R^2$; the coalition argument is unchanged. In the
$i$-th swap, the first-order terms give
$\sum_j\E[\partial_jf(W_i)]\,(\E X_{i,j} - \E N_{i,j})/
\sqrt n = 0$ and the second-order terms give
$\frac1{2n}\sum_{j,k}\E[\partial_{jk}f(W_i)]\,(\Sigma_{jk}
- \Sigma_{jk}) = 0$: means and covariances match. Question
15 bounds the two remainders:
\[
\bigl|\E f(H_i) - \E f(H_{i-1})\bigr| \leq
\frac{\sqrt2M_3}3\cdot\frac{\E\norm{X_i}^3 +
\E\norm{N_i}^3}{n^{3/2}}
= \frac{\sqrt2M_3(\beta' + \gamma')}{3\,n^{3/2}},
\]
and telescoping over the $n$ swaps:
\[
\Bigl|\E f\Bigl(\frac{S_n}{\sqrt n}\Bigr) - \E
f(G_n)\Bigr| \leq \frac{\sqrt2\,M_3(\beta' +
\gamma')}{3\sqrt n},
\qquad G_n \sim \mathcal N(0, \Sigma)\ \text{exactly} .
\]
Upgrade: $\E\norm{T_n}^2 = \E\norm{X_1}^2 =
\operatorname{tr}\Sigma$ (cross terms vanish by
\omterm{def:b3:probability:independence}{independence} and centering), so $\P(\norm{T_n} > A) \leq
\operatorname{tr}\Sigma/A^2$, and likewise for $N$:
tightness. Given a bounded \omterm{def:b3:topology:continuity}{continuous} $f$ and $\varepsilon
> 0$, multiply by a smooth plateau $\chi$ equal to $1$ on
the ball of radius $A$ and supported in radius $A + 1$
(mollify an indicator in $\R^2$,
\cref{thm:b3:lp:regularization}); $g = f\chi$ is uniformly
\omterm{def:b3:topology:continuity}{continuous} with \omterm{def:b3:topology:compact}{compact} support, so its two-dimensional
mollification $g_\eta$ is $\mathcal C^\infty$ with bounded
derivatives of all orders and $\norm{g - g_\eta}_\infty
\leq \varepsilon$ for small $\eta$. The three-$\varepsilon$
chain of question 5 then transfers verbatim: $\E f(T_n)
\to \E f(N)$ for every bounded \omterm{def:b3:topology:continuity}{continuous} $f \colon \R^2
\to \R$. This is \cref{thm:b3:clt:multiclt} for $d = 2$,
now proved --- swapping sidesteps the two-dimensional
L\'evy theorem that the chapter had left admitted.

\textbf{17.} For bounded \omterm{def:b3:topology:continuity}{continuous} $g \colon \R \to \R$,
the map $x \mapsto g(\langle t, x\rangle)$ is bounded
\omterm{def:b3:topology:continuity}{continuous} on $\R^2$, so question 16 gives $\E g(\langle
t, T_n\rangle) \to \E g(\langle t, N\rangle)$: every
projection converges in distribution, and $\langle t,
N\rangle \sim \mathcal N(0, t^{\mathsf T}\Sigma t)$. (This
is the easy direction of Cram\'er--Wold: joint convergence
implies convergence of all linear images.) Application:
$V_i = (\xi_i, \xi_i^2 - 1)$ are i.i.d.\ centered vectors
($\E\xi_1^2 = 1$), with covariance entries $\V(\xi_1) =
1$, $\operatorname{Cov}(\xi_1, \xi_1^2 - 1) = \E\xi_1^3$
and $\V(\xi_1^2 - 1) = \E\xi_1^4 - 1$; the third moment
$\E\norm{V_1}^3 \leq 4\bigl(\E\abs{\xi_1}^3 +
\E\abs{\xi_1^2 - 1}^3\bigr)$ is finite when $\xi_1 \in
L^6$. Question 16 yields the displayed joint \omterm{def:b3:clt:gaussianvector}{Gaussian}
limit, and \cref{thm:b3:clt:gaussianvector}(2): the two
limit coordinates are \omterm{def:b3:probability:independence}{independent} exactly when the
covariance $\E\xi_1^3$ vanishes --- for symmetric laws,
empirical mean and empirical variance decouple
asymptotically.

\textbf{18.} Write $g(x) - g(\theta) = (g'(\theta) +
\eta(x))(x - \theta)$ where $\eta(x) = \frac{g(x) -
g(\theta)}{x - \theta} - g'(\theta)$ for $x \neq \theta$
and $\eta(\theta) = 0$: differentiability at $\theta$
means precisely $\eta(x) \to 0$ as $x \to \theta$.
\emph{Step 1:} $\hat\theta_n \to \theta$ in probability:
for $\varepsilon > 0$ and any $A > 0$, eventually
$\varepsilon\sqrt n \geq A$, so $\P(\abs{\hat\theta_n -
\theta} > \varepsilon) \leq \P(\abs{\sqrt
n(\hat\theta_n - \theta)} > A) \to \P(\sigma\abs N > A)$
(distribution functions converge at the \omterm{def:b3:topology:continuity}{continuity} points
$\pm A$), and the right side tends to $0$ as $A \to
\infty$. \emph{Step 2:} $\eta(\hat\theta_n) \to 0$ in
probability: given $\varepsilon' > 0$, pick $\delta$ with
$\abs\eta \leq \varepsilon'$ on $\abs{x - \theta} \leq
\delta$; then $\P(\abs{\eta(\hat\theta_n)} > \varepsilon')
\leq \P(\abs{\hat\theta_n - \theta} > \delta) \to 0$.
\emph{Step 3:}
\[
\sqrt n\bigl(g(\hat\theta_n) - g(\theta)\bigr) =
g'(\theta)\,\sqrt n(\hat\theta_n - \theta) +
\eta(\hat\theta_n)\cdot\sqrt n(\hat\theta_n - \theta) .
\]
By Slutsky's product rule (\cref{exo:b3:clt:8}, with the
sequence $\eta(\hat\theta_n) \to 0$ in probability and the
convergent-in-law $\sqrt n(\hat\theta_n - \theta)$), the
second term converges in law to $0\cdot\mathcal N(0,
\sigma^2) = 0$, hence to $0$ in probability
(\cref{exo:b3:clt:4}(b)); the first converges in law to
$g'(\theta)\mathcal N(0, \sigma^2)$ (Slutsky again, or the
affine rule for \omterm{def:b3:clt:cf}{characteristic functions}); Slutsky's sum
rule assembles them: the limit is $\mathcal N(0,
g'(\theta)^2\sigma^2)$.

\textbf{19.} (a) The CLT gives $\sqrt n(\bar X_n - \mu)
\Rightarrow \mathcal N(0, \sigma^2)$; the delta method
with $g(x) = x^2$, $g'(\mu) = 2\mu$, gives $\sqrt n(\bar
X_n^2 - \mu^2) \Rightarrow \mathcal N(0, 4\mu^2\sigma^2)$
--- degenerate (limit $0$) when $\mu = 0$. In that case
the fluctuation lives one scale up: $n\bar X_n^2 =
(\sqrt n\,\bar X_n)^2$, and for $t > 0$
\[
\P\bigl(n\bar X_n^2 \leq t\bigr) = \P\bigl(-\sqrt t \leq
\sqrt n\,\bar X_n \leq \sqrt t\bigr) \longrightarrow
\Phi\Bigl(\frac{\sqrt t}\sigma\Bigr) -
\Phi\Bigl(-\frac{\sqrt t}\sigma\Bigr) = \P(\sigma^2N^2
\leq t) :
\]
$n\bar X_n^2 \Rightarrow \sigma^2N^2$, the square of a
\omterm{def:b3:clt:gaussianvector}{Gaussian} (a ``chi-squared'' law) --- when the first
derivative dies, the second-order term of Taylor dictates
a \omterm{def:b3:clt:gaussianvector}{non-Gaussian} limit. (b) Here $\sqrt n(\hat p_n - p)
\Rightarrow \mathcal N(0, p(1 - p))$ and $g(p) =
\arcsin\sqrt p$ has $g'(p) = \frac1{2\sqrt{p(1 - p)}}$, so
$g'(p)^2\,p(1 - p) = \frac14$: the limit is $\mathcal
N(0, \frac14)$ for every $p \in \intoo01$. On the
$\arcsin$ scale the asymptotic $95\%$ error bar is $\pm
\frac{0.98}{\sqrt n}$, known in advance --- whereas in
\cref{ex:b3:clt:confidence} the width involved the unknown
$\sigma = \sqrt{p(1-p)}$, to be worst-cased by $\frac12$
or estimated: the transformation \emph{stabilizes} the
variance.

\textbf{20.} Let $A^* = \{k : \mu(\{k\}) > \nu(\{k\})\}$
and $\Delta_k = \mu(\{k\}) - \nu(\{k\})$, so
$\sum_k\Delta_k = 0$. For any $A \subseteq \N$: $\mu(A) -
\nu(A) = \sum_{k\in A}\Delta_k \leq \sum_{k\in
A^*}\Delta_k$, with equality at $A = A^*$; and since the
positive and negative parts of $(\Delta_k)$ have equal
total mass, $\sum_{A^*}\Delta_k =
\frac12\sum_k\abs{\Delta_k}$. Exchanging $\mu, \nu$ handles
the sign: $d_{\mathrm{TV}}(\mu, \nu) =
\frac12\sum_k\abs{\Delta_k}$. Coupling: for any $A$,
\[
\mu(A) - \nu(A) = \E\bigl[\mathbf 1_A(X) - \mathbf
1_A(Y)\bigr] = \E\bigl[(\mathbf 1_A(X) - \mathbf
1_A(Y))\,\mathbf 1_{X\neq Y}\bigr] \leq \P(X \neq Y),
\]
and take the supremum over $A$.

\textbf{21.} The two laws charge: $k = 0$: $1 - p$ versus
$\eu^{-p}$, with $\eu^{-p} > 1 - p$; $k = 1$: $p$ versus
$p\,\eu^{-p} < p$; $k \geq 2$: $0$ versus the Poisson
remainder $1 - \eu^{-p} - p\eu^{-p} \geq 0$. Hence
\[
d_{\mathrm{TV}} = \tfrac12\bigl[(\eu^{-p} - 1 + p) + (p -
p\eu^{-p}) + (1 - \eu^{-p} - p\eu^{-p})\bigr] =
\tfrac12\bigl(2p - 2p\eu^{-p}\bigr) = p(1 - \eu^{-p}),
\]
and $1 - \eu^{-p} \leq p$ gives the bound $p^2$.

\textbf{22.} Write $H_{i-1} = W_i + X_i$ and $H_i = W_i +
Y_i$ with $W_i = \sum_{j<i}Y_j + \sum_{j>i}X_j$,
\omterm{def:b3:probability:independence}{independent} of the pair $(X_i, Y_i)$ (coalitions). For $A
\subseteq \N$, conditioning on the countably many values
by \omterm{def:b3:probability:independence}{independence},
\[
\P(H_{i-1} \in A) = \sum_{k\geq0}\P(X_i = k)\,\P(W_i + k
\in A),
\]
and likewise for $H_i$ with $Y_i$. Subtracting, with $c_k
= \P(W_i + k \in A) \in \intcc01$ and $\Delta_k = \P(X_i =
k) - \P(Y_i = k)$ of zero sum:
\[
\abs{\P(H_{i-1} \in A) - \P(H_i \in A)} =
\Bigl|\sum_k\Delta_k\bigl(c_k - \tfrac12\bigr)\Bigr| \leq
\tfrac12\sum_k\abs{\Delta_k} =
d_{\mathrm{TV}}\bigl(\mathcal B(1, p_i), \mathcal
P(p_i)\bigr) .
\]
Telescoping from $H_0 = S$ to $H_n = \sum_iY_i \sim
\mathcal P(\lambda)$ (\cref{exo:b3:clt:1}, iterated) and
using question 21:
\[
\abs{\P(S \in A) - \P(\mathcal P(\lambda) \in A)} \leq
\sum_{i=1}^np_i\bigl(1 - \eu^{-p_i}\bigr) \leq
\sum_{i=1}^np_i^2
\qquad\text{for every } A :
\]
Le Cam's inequality. (Question 20's coupling bound gives an
alternative route: couple each pair on one uniform
variable so that $\P(X_i \neq Y_i) \leq p_i^2$ and bound
$\P(S \neq \sum Y_i)$; swapping needs no construction at
all.)

\textbf{23.} (a) With $p_i = \frac\lambda n$:
$d_{\mathrm{TV}}(\text{law of }S, \mathcal P(\lambda))
\leq \frac{\lambda^2}n$. This sharpens
\cref{exo:b3:clt:5} three times over: an explicit error at
every finite $n$, uniformity over all events $A$ at once
(not one interval at a time), and no need for equal
$p_i$ --- only $\sum_ip_i^2$ small, e.g.\ $\sum p_i^2 \leq
\lambda\max_ip_i$: \emph{many rare events, none
dominant}. (b) Here $n = 500$, $p_i = \frac1{500}$,
$\lambda = 1$: the Poisson model errs by at most
$500\cdot\frac1{500^2} = 0.002$ on every event; in
particular, taking $A = \{0\}$,
\[
\P(\text{no letter astray}) = \Bigl(1 -
\frac1{500}\Bigr)^{500},
\qquad
\Bigl|\P(\text{no letter astray}) - \eu^{-1}\Bigr| \leq
0.002,
\]
so the answer is $\eu^{-1} \approx 0.368$ up to a
guaranteed $0.002$ (the true discrepancy is about
$4\cdot10^{-4}$). (c) The
problem closes on one method with two regimes. When $n$
comparable contributions each carry variance
$\frac1n$, matching \emph{two} moments against the
\omterm{def:b3:clt:gaussianvector}{Gaussian} makes the swap errors $o(\frac1n)$ each: sums go
\omterm{def:b3:clt:gaussianvector}{Gaussian} --- with Taylor as the local comparison tool.
When $n$ contributions are indicators of probability
$p_i$, matching the \emph{mean} against a Poisson atom
makes each swap cost $p_i^2$: counts of rare events go
Poisson --- with total variation as the exact local
comparison. Same hybrids, same telescope, different local
estimate: replacement is a strategy, not a theorem, and
the \omterm{def:b3:clt:gaussianvector}{Gaussian} and Poisson limits are its two oldest
dividends.

\textbf{24.} The CLT gives $\sqrt n(\bar X_n - \mu)
\Rightarrow \mathcal N(0, \sigma^2)$, and $g(x) = \ln x$ is
differentiable at $\mu > 0$ with $g'(\mu) = \frac1\mu$: the
delta method (Part VI) yields $\sqrt n(\ln\bar X_n - \ln\mu)
\Rightarrow \mathcal N(0, \sigma^2/\mu^2)$. Unwinding the
interval $\abs{\ln\bar X_n - \ln\mu} \leq
\frac{1.96\,\sigma}{\mu\sqrt n}$ by exponentiation:
\[
\mu \in \bar X_n\cdot
\eu^{\pm1.96\,\sigma/(\mu\sqrt n)}
\qquad\text{with asymptotic probability } 95\%
\]
(in practice $\sigma/\mu$ is replaced by its empirical
version, Slutsky as in \cref{exo:b3:clt:8}). The
multiplicative interval is the natural one when the data
are positive with errors proportional to their size ---
incomes, concentrations, half-lives: quantities that live
on a log scale, where symmetric additive intervals could
even cross zero.

\textbf{25.} $\E[(X - p)^3] = (1-p)^3p + (-p)^3(1 - p) =
p(1-p)\bigl[(1-p)^2 - p^2\bigr] = p(1-p)(1 - 2p)$. In the
swapping analysis (Part IV), the leading error term after
matching two moments carries the \emph{signed} third
moment: for $p < \frac12$ it is positive (the law leans
right: rare large excursions above the mean), and the
normal approximation systematically misplaces mass ---
underestimating the short left tail and overestimating the
right --- with an error of order $n^{-1/2}$; at $p =
\frac12$ the third moment vanishes, the Bernoulli matches
the \omterm{def:b3:clt:gaussianvector}{Gaussian} to third order, and the rate improves (Part
IV's matched-moment question). Numerically: $\P(S = 0) =
0.9^{20} = 0.1216$, while the \omterm{def:b3:clt:gaussianvector}{Gaussian}
$\mathcal N(2, 1.8)$ gives $\Phi\bigl(\frac{0.5 -
2}{\sqrt{1.8}}\bigr) = \Phi(-1.118) \approx 0.132$: the
normal curve, ignorant of the wall at $0$ and of the
rightward skew, puts too much mass at the bottom --- the
predicted sign of the error, visible at $n = 20$.
\end{solution}
